I follow the hint. The only standard convex-geometric input I use is that every vertex of a polyhedron is an exposed point, i.e. it is the unique maximizer of some linear functional; see Proposition \(3.3\) in the cited notes. ([ms.uky.edu](https://www.ms.uky.edu/~Lee/ma714fa13/notes.pdf))

For each \(s\in S\), choose a linear functional \(\phi_s:\mathbb R^2\to\mathbb R\) that is uniquely maximized at \(s\) on \(P=\operatorname{conv}(S)\), hence also on \(S\). Since \(T\) is finite, choose \(t_s\in T\) with
\[
\phi_s(t_s)=\max_{t\in T}\phi_s(t),
\]
and define
\[
\ell_s:=s+t_s.
\]

Now prove the six assertions.

---

### \(1.\) Lonely points exist

Fix \(s\in S\). I claim that \(\ell_s\) is lonely.

Indeed, suppose
\[
\ell_s=s'+t'
\qquad\text{for some }s'\in S,\ t'\in T.
\]
Applying \(\phi_s\), we get
\[
\phi_s(s)+\phi_s(t_s)
=\phi_s(\ell_s)
=\phi_s(s')+\phi_s(t').
\]
But by construction,
\[
\phi_s(s')\le \phi_s(s)
\quad\text{and}\quad
\phi_s(t')\le \phi_s(t_s).
\]
Hence equality must hold in both inequalities. Since \(s\) is the **unique** maximizer of \(\phi_s\) on \(S\), we obtain
\[
s'=s.
\]
Therefore
\[
t'=\ell_s-s=t_s.
\]
So there is a unique \(t\in T\) with \(\ell_s-t\in S\), namely \(t=t_s\). Thus \(\ell_s\) is lonely.

Notice that this did **not** require \(t_s\) to be a unique maximizer of \(\phi_s\) on \(T\).

---

### \(2.\) Distinctness

Suppose
\[
\ell_{s_1}=\ell_{s_2}.
\]
Then this common point has two representations:
\[
\ell_{s_1}=s_1+t_{s_1}=s_2+t_{s_2}.
\]
By part \(1\), the point \(\ell_{s_1}\) is lonely, so its translation vector is unique. Hence
\[
t_{s_1}=t_{s_2}.
\]
Subtracting, we get
\[
s_1=s_2.
\]
Therefore the points \(\ell_s\) are pairwise distinct.

Since \(|S|=n\), the set
\[
L:=\{\ell_s:s\in S\}
\]
has exactly \(n\) elements.

---

### \(3.\) Non-vanishing

By part \(1\), the only \(t\in T\) for which \(\ell_s-t\in S\) is \(t=t_s\). Therefore
\[
F(\ell_s)
=
\sum_{\substack{t\in T\\ \ell_s-t\in S}} \varepsilon(t)
=
\varepsilon(t_s).
\]
Since \(\varepsilon(t_s)\in\{+1,-1\}\), we have
\[
F(\ell_s)=\varepsilon(t_s)\neq 0.
\]

---

### \(4.\) Lower bound

Let
\[
N:=\{p\in\mathbb R^2:F(p)\neq 0\}.
\]
By part \(3\), every \(\ell_s\) lies in \(N\). By part \(2\), the points \(\ell_s\) are pairwise distinct. Hence
\[
|N|\ge |L|=n.
\]

This is exactly the desired lower bound:
\[
|\{p\in\mathbb R^2:F(p)\neq 0\}|\ge n.
\]

---

### \(5.\) Rigidity in the equality case

Assume now that
\[
|N|=n.
\]
From part \(4\), we already know that
\[
L\subseteq N
\quad\text{and}\quad
|L|=n.
\]
Since also \(|N|=n\), it follows that
\[
N=L.
\]
So
\[
\{p\in\mathbb R^2:F(p)\neq 0\}
=
\{\ell_s:s\in S\}.
\]

In particular, if \(p\in S+T\) is not of the form \(\ell_s\), then \(p\notin N\), hence
\[
F(p)=0.
\]

---

### \(6.\) Vertices survive

In fact, I will prove the stronger statement that **every vertex of \(P+Q\) already satisfies \(F(v)\neq 0\)**, even without assuming \(|N|=n\). The stated conclusion then follows immediately.

Let \(v\) be a vertex of \(P+Q\). Again using the exposed-point fact for vertices, choose a linear functional \(\psi:\mathbb R^2\to\mathbb R\) that is uniquely maximized at \(v\) on \(P+Q\). ([ms.uky.edu](https://www.ms.uky.edu/~Lee/ma714fa13/notes.pdf))

Set
\[
M_S:=\max_{s\in S}\psi(s),
\qquad
M_T:=\max_{t\in T}\psi(t),
\]
and choose \(s^\ast\in S\), \(t^\ast\in T\) such that
\[
\psi(s^\ast)=M_S,
\qquad
\psi(t^\ast)=M_T.
\]

Now take any \(x\in P+Q\). Then \(x=p+q\) for some \(p\in P\), \(q\in Q\). Since \(p\in\operatorname{conv}(S)\), \(p\) is a convex combination of points of \(S\), so
\[
\psi(p)\le M_S.
\]
Similarly,
\[
\psi(q)\le M_T.
\]
Therefore
\[
\psi(x)=\psi(p)+\psi(q)\le M_S+M_T
=\psi(s^\ast+t^\ast).
\]
Thus \(s^\ast+t^\ast\) is a maximizer of \(\psi\) on \(P+Q\). By uniqueness of the maximizer,
\[
v=s^\ast+t^\ast.
\]

Next I claim that \(s^\ast\) is the unique maximizer of \(\psi\) on \(S\). If there were some \(s'\neq s^\ast\) with
\[
\psi(s')=M_S,
\]
then
\[
\psi(s'+t^\ast)=M_S+M_T=\psi(v),
\]
so \(s'+t^\ast\) would also be a maximizer of \(\psi\) on \(P+Q\). But
\[
s'+t^\ast\neq s^\ast+t^\ast=v,
\]
contradicting uniqueness. Hence \(s^\ast\) is unique. The same argument shows that \(t^\ast\) is the unique maximizer of \(\psi\) on \(T\).

Now suppose
\[
v=s'+t'
\qquad\text{with }s'\in S,\ t'\in T.
\]
Then
\[
\psi(v)=\psi(s')+\psi(t')\le M_S+M_T=\psi(v),
\]
so equality holds in both inequalities. Hence
\[
\psi(s')=M_S
\quad\text{and}\quad
\psi(t')=M_T.
\]
By uniqueness of the maximizers on \(S\) and \(T\), this forces
\[
s'=s^\ast,
\qquad
t'=t^\ast.
\]
Thus \(v\) is lonely. Consequently,
\[
F(v)=\varepsilon(t^\ast)\neq 0.
\]

So every vertex \(v\) of \(P+Q\) lies in \(N\). Under the hypothesis
\[
|N|=n,
\]
part \(5\) gives
\[
N=\{\ell_s:s\in S\},
\]
hence every vertex of \(P+Q\) belongs to \(\{p:F(p)\neq 0\}\), as required.

---

Thus all six statements are proved.